% =========================================================
% Project 1 — Atlanta temperatures: full statistical study
% Time Series Analysis — Chapter 1
% =========================================================
\documentclass[a4paper,11pt]{article}

\newcommand{\TPnumero}{Project 1}
\newcommand{\macrosPath}{../../preamble/macros.tex}
% =========================================================
% Tutorial / lab preamble — Time Series Analysis module
% article class + fancyhdr + tcolorbox + listings (English)
% Author : Aymen Ben Brik (aymen.benbrik@esprit.tn)
% =========================================================

\usepackage[utf8]{inputenc}
\usepackage[T1]{fontenc}
\usepackage[english]{babel}
\usepackage{lmodern}
\usepackage{amsmath, amssymb, amsthm, mathtools, bm}
\usepackage{graphicx}
\usepackage{booktabs}
\usepackage{array}
\usepackage{multirow}
\usepackage{colortbl}
\usepackage{xcolor}
\usepackage{tikz}
\usetikzlibrary{positioning, arrows.meta, calc, shapes, fit, backgrounds, matrix}
\usepackage{listings}
\usepackage[most]{tcolorbox}
\usepackage{geometry}
\geometry{a4paper, margin=2cm, headheight=15pt}
\usepackage{fancyhdr}
\usepackage[hidelinks]{hyperref}
\usepackage{enumitem}

% --- Colours ---
\definecolor{espritBlue}{RGB}{30, 60, 110}
\definecolor{espritAccent}{RGB}{200, 80, 30}
\definecolor{boxEx}{RGB}{30, 60, 110}
\definecolor{boxQ}{RGB}{180, 120, 0}
\definecolor{boxC}{RGB}{30, 100, 60}
\definecolor{boxAtt}{RGB}{200, 80, 30}

% --- Listings ---
\definecolor{codebg}{RGB}{248,248,248}
\definecolor{codeframe}{RGB}{220,220,220}
\definecolor{keyword}{RGB}{0,82,155}
\definecolor{comment}{RGB}{88,110,117}
\definecolor{string}{RGB}{133,153,0}

\lstdefinestyle{pythonstyle}{
  language=Python,
  basicstyle=\ttfamily\small,
  keywordstyle=\color{keyword}\bfseries,
  commentstyle=\color{comment}\itshape,
  stringstyle=\color{string},
  backgroundcolor=\color{codebg},
  frame=single,
  rulecolor=\color{codeframe},
  frameround=tttt,
  breaklines=true,
  showstringspaces=false,
  tabsize=2
}
\lstdefinestyle{Rstyle}{
  language=R,
  basicstyle=\ttfamily\small,
  keywordstyle=\color{keyword}\bfseries,
  commentstyle=\color{comment}\itshape,
  stringstyle=\color{string},
  backgroundcolor=\color{codebg},
  frame=single,
  rulecolor=\color{codeframe},
  frameround=tttt,
  breaklines=true,
  showstringspaces=false,
  tabsize=2
}
\lstset{style=pythonstyle}

% --- Common macros (configurable path) ---
\providecommand{\macrosPath}{../../preamble/macros.tex}
\input{\macrosPath}

% --- Exercise counter ---
\newcounter{excounter}
\renewcommand{\theexcounter}{\arabic{excounter}}

\newtcolorbox[use counter=excounter]{exercise}[1][]{
  enhanced, breakable,
  colback=boxEx!5, colframe=boxEx, fonttitle=\bfseries,
  title={Exercise~\theexcounter\ifstrempty{#1}{}{ -- #1}},
  arc=2pt, boxrule=0.6pt
}
\newtcolorbox{question}[1][]{
  enhanced, breakable,
  colback=boxQ!5, colframe=boxQ, fonttitle=\bfseries,
  title={Question\ifstrempty{#1}{}{ -- #1}},
  arc=2pt, boxrule=0.5pt
}
\newtcolorbox{solution}[1][]{
  enhanced, breakable,
  colback=boxC!5, colframe=boxC, fonttitle=\bfseries,
  title={Solution\ifstrempty{#1}{}{ -- #1}},
  arc=2pt, boxrule=0.5pt
}
\newtcolorbox{warning}[1][]{
  enhanced, breakable,
  colback=boxAtt!5, colframe=boxAtt, fonttitle=\bfseries,
  title={Warning\ifstrempty{#1}{}{ -- #1}},
  arc=2pt, boxrule=0.5pt
}

% --- Headers / footers ---
\pagestyle{fancy}
\fancyhf{}
\fancyhead[L]{\textsc{\moduleNom} -- \auteurNom}
\fancyhead[R]{\TPnumero}
\fancyfoot[C]{\thepage}
\fancyfoot[L]{\auteurInstitution}
\fancyfoot[R]{\auteurEmail}
\renewcommand{\headrulewidth}{0.4pt}
\renewcommand{\footrulewidth}{0.2pt}

% --- Placeholder ---
\providecommand{\TPnumero}{Lab}


\title{Project 1: Statistical study of Atlanta monthly temperatures\\
\large Descriptive statistics, normality, inference}
\author{\auteurNom\\\auteurEmail\\\auteurInstitution}
\date{\anneeUniversitaire}

\begin{document}
\maketitle

\section*{Context}

The dataset \texttt{AvTempAtlanta.txt} contains the monthly average
temperatures (degrees Fahrenheit) of Atlanta, USA, over several
decades. The goal of this project is to apply the basic statistical
tools of Chapter~1 to a real dataset: descriptive statistics, normality
diagnostic, confidence intervals, and hypothesis testing.

\section*{Deliverables}
\begin{itemize}
  \item A reproducible Python script (\texttt{project1\_<name>.py})
        with seed $42$.
  \item A 4--6 page PDF report including all figures, tables and
        discussion.
  \item Optional 5-minute oral defence per student.
\end{itemize}

% =========================================================
\begin{exercise}[Descriptive statistics]
\textbf{(1)} Load the data, plot the time series. Comment on the
qualitative features (trend, seasonality, range).

\textbf{(2)} Compute and report:
\begin{itemize}
  \item sample mean, variance, standard deviation, min, max, median;
  \item sample skewness and kurtosis.
\end{itemize}
Plot a histogram with kernel density estimate overlay.
\end{exercise}

% =========================================================
\begin{exercise}[Normality diagnostic]
\textbf{(1)} Plot a Q--Q plot (\texttt{scipy.stats.probplot}) against
the standard normal.

\textbf{(2)} Run two formal normality tests:
\begin{itemize}
  \item Shapiro--Wilk (\texttt{scipy.stats.shapiro});
  \item Jarque--Bera (\texttt{scipy.stats.jarque\_bera}).
\end{itemize}
Report the test statistics and $p$-values. Conclude on the normality
assumption at the $5\%$ level.
\end{exercise}

% =========================================================
\begin{exercise}[Confidence interval and hypothesis test]
\textbf{(1)} Compute the 90\%, 95\% and 99\% $t$-confidence intervals
for the mean monthly temperature.

\textbf{(2)} Test $H_0 : \mu = 60$ vs.\ $H_1 : \mu \neq 60$ at level
$5\%$. Report $T$, the critical value, the $p$-value, and conclude.

\textbf{(3)} Repeat with $H_0 : \mu = 65$. Discuss the role of $\mu_0$
in the test.
\end{exercise}

% =========================================================
\begin{exercise}[Yearly aggregation]
\textbf{(1)} Aggregate the monthly data into yearly averages (12-month
mean). How many yearly observations do you obtain?

\textbf{(2)} Run a $t$-test on the yearly averages: $H_0 : \mu = 60$
vs.\ $H_1 : \mu \neq 60$. How does the conclusion compare with the
monthly version?

\textbf{(3)} Why does aggregation change the test sensitivity?
Discuss in terms of the variance of the sample mean.
\end{exercise}

% =========================================================
\begin{exercise}[Bonus: bootstrap and assumption robustness]
\textbf{(1)} Build a bootstrap CI for the mean ($B = 10^4$). Compare
with the $t$-CI.

\textbf{(2)} If the data are not normal (cf.\ Exercise~2), is the
$t$-CI still valid? Justify (CLT, sample size).
\end{exercise}

% =========================================================
\section*{Marking grid (out of 100)}

\begin{center}
\begin{tabular}{p{8cm}rl}
\toprule
\textbf{Criterion} & \textbf{Points} & \\
\midrule
Loading + descriptive statistics + plots               & 15 \\
Normality diagnostic + interpretation                  & 15 \\
$t$-CI (3 levels) + $t$-test ($\mu_0 \in \{60, 65\}$)  & 25 \\
Yearly aggregation + comparison                        & 15 \\
Bootstrap (bonus) + robustness discussion              & 10 \\
Reproducibility (seed, code structure, comments)       & 10 \\
Report quality (figures, tables, discussion)           & 10 \\
\midrule
\textbf{Total}                                         & \textbf{100} \\
\bottomrule
\end{tabular}
\end{center}

\bigskip
\hfill\auteurNom\par
\hfill\auteurEmail\par

\end{document}
